\uuid{DTGM}
\titre{Domaine de définition}
\auteur{bodin ; hanani ; mzari}
\organisation{exo7}
\licence{CC BY-NC-SA 4.0}
\datecreate{2020-01-20}
\niveau{L1}
\module{Analyse}
\chapitre{Fonctions circulaires et hyperboliques inverses}
\sousChapitre{Autre}
\difficulte{1}
\theme{Fonctions usuelles, domaine de définition}
\exo7qcmid{lille-1.fonctions-usuelles.008}

\contenu{
\begin{question}
\qclassification{126.00}
\qtags{motcle=domaine de définition}
Soit $(E)$ l'équation : $e^{2x}+e^x-2=0$. Quelles sont les assertions vraies ?

\begin{answers}
    \good{$(E)$ est définie sur $\Rr$.}

    \bad{Le domaine de définition de $(E)$ est $\Rr^+$.}

    \bad{$(E)$ admet deux solutions distinctes.}

    \good{$(E)$ admet une unique solution $x=0$.}
\end{answers}
\begin{explanations}
La fonction exponentielle est définie sur $\Rr$, donc $(E)$ est définie sur $\Rr$.
En posant $y=e^x$, on se ramène à résoudre l'équation du second degré $y^2+y-2=0$. En résolvant cette équation, on obtient $y=1$ ou $y=-2$. Par conséquent, $x=0$.
\end{explanations}
\end{question}
}
